% Part: lambda-calculus
% Chapter: syntax
% Section: beta

\documentclass[../../../include/open-logic-section]{subfiles}

\begin{document}

\olfileid{lam}{syn}{bet}
\olsection{کاهش~$\beta$}

با دیدن~$(\lambd[m][(\lambd[y][y]) m])$، طبیعی است حدس بزنیم که این
ترم با~$\lambd[m][m]$ رابطه‌ای دارد؛ یعنی ترم دوم باید حاصل «ساده‌سازی»
ترم نخست باشد. مفهوم \emph{کاهش~$\beta$} این شهود را به‌طور صوری بیان
می‌کند.

\begin{defn}[انقباض~$\beta$، $\bredone$] \ollabel{defn:betacontr}
  \emph{انقباض~$\beta$} ($\bredone$) کوچک‌ترین رابطهٔ سازگار روی ترم‌هاست
  که شرط زیر را برآورده می‌کند:
  \[
    (\lambd[x][N])Q \bredone \Subst{N}{Q}{x} 
  \]
  می‌گوییم~$P$ به~$Q$ \emph{منقبض به~$\beta$} می‌شود اگر~$P \bredone Q$.
  ترمی به‌صورت~$(\lambd[x][N])Q$ یک \emph{سرخ‌عبارت} نامیده می‌شود.
\end{defn}

\begin{prob} \ollabel{prob:def}
  تعریف استقرایی هم‌ارز انقباض~$\beta$ را، همان‌گونه که برای تغییر متغیر
  مقید در~\olref[lam][syn][alp]{defn:aconvone} انجام دادیم، به‌تفصیل
  بنویسید.
\end{prob}
  
\begin{defn}[کاهش~$\beta$، $\bred$] \ollabel{defn:betared}
  \emph{کاهش~$\beta$} ($\bred$) کوچک‌ترین رابطهٔ بازتابی و تعدی روی
  ترم‌هاست که~$\bredone$ را در بر می‌گیرد. می‌گوییم~$P$ به~$Q$
  کاهش~$\beta$ می‌یابد اگر~$P \bred Q$.
\end{defn}

هرگاه بافت روشن باشد، $\redone$ را به‌جای~$\bredone$، و~$\red$ را
به‌جای~$\bred$ می‌نویسیم.

به‌طور غیرصوری، $M \bred N$ اگر و تنها اگر بتوان~$M$ را با صفر یا چند
گام انقباض~$\beta$ به~$N$ تبدیل کرد.

\begin{defn}[نرمال نسبت به~$\beta$]
ترمی که نتوان آن را بیش از این منقبض به~$\beta$ کرد
\emph{نرمال نسبت به~$\beta$} نامیده می‌شود.
\end{defn}

اگر~$M \bred N$ و~$N$ نرمال نسبت به~$\beta$ باشد، می‌گوییم~$N$ یک
\emph{صورت نرمال}~$M$ است. ممکن است بپرسیم آیا صورت نرمال یک ترم یکتا
است؛ پاسخ مثبت است، چنان‌که بعدتر خواهیم دید.

چند مثال را در نظر می‌گیریم.
\begin{enumerate}
\item داریم
\begin{align*}
(\lambd[x][xxy]) \lambd[z][z] & \redone (\lambd[z][z])(\lambd[z][z]) y \\
& \redone (\lambd[z][z]) y \\
& \redone y
\end{align*}
\item «ساده‌سازی» یک ترم ممکن است در واقع آن را پیچیده‌تر کند:
\begin{align*}
(\lambd[x][xxy])(\lambd[x][xxy]) & \redone (\lambd[x][xxy])(\lambd[x][xxy])y \\
& \redone (\lambd[x][xxy])(\lambd[x][xxy])yy \\
& \redone \dots
\end{align*}
\item همچنین ممکن است ترمی را بی‌تغییر بگذارد:
\[
(\lambd[x][xx])(\lambd[x][xx]) \redone (\lambd[x][xx])(\lambd[x][xx])
\]
\item افزون بر این، برخی ترم‌ها را می‌توان به بیش از یک شیوه کاهش داد؛
  برای مثال،
\[
(\lambd[x][(\lambd[y][yx]) z]) v \redone (\lambd[y][yv]) z
\]
با منقبض‌کردن بیرونی‌ترین اعمال؛ و
\[
(\lambd[x][(\lambd[y][yx]) z]) v \redone (\lambd[x][zx]) v
\]
با منقبض‌کردن درونی‌ترین اعمال. بااین‌حال توجه کنید که در این حالت هر
دو ترم سپس به ترم یکسان~$zv$ کاهش می‌یابند.
\end{enumerate}

حاصل نهایی در مثال آخر تصادفی نیست، بلکه خاصیتی ژرف و مهم از حساب
لامبدا را نشان می‌دهد که خاصیت چرچ--راسر نام دارد.

\begin{digress}
  به‌طور کلی بیش از یک راه برای کاهش~$\beta$ یک ترم وجود دارد؛ ازاین‌رو
  راهبردهای کاهش فراوانی ابداع شده‌اند که رایج‌ترین آن‌ها
  \emph{راهبرد طبیعی} است. راهبرد طبیعی همواره \emph{چپ‌ترین}
  سرخ‌عبارت را منقبض می‌کند؛ جایگاه سرخ‌عبارت نقطهٔ آغاز آن در ترم تعریف
  می‌شود. راهبرد طبیعی این خاصیت سودمند را دارد که یک ترم با راهبردی
  به صورت نرمال کاهش‌پذیر است اگر و تنها اگر با راهبرد طبیعی به صورت
  نرمال کاهش‌پذیر باشد. در ادامه، مگر آنکه خلافش تصریح شود، راهبرد طبیعی
  را به کار خواهیم برد.
\end{digress}

\begin{defn}[هم‌ارزی~$\beta$، $\equal$]
  \emph{هم‌ارزی~$\beta$} ($\equal$) رابطه‌ای است که به‌صورت استقرایی
  چنین تعریف می‌شود:
  \begin{enumerate}
  \item $M \equal M$.
  \item اگر~$M \equal N$، آنگاه~$N \equal M$.
  \item اگر~$M \equal N$ و~$N \equal O$، آنگاه~$M \equal O$.
  \item اگر~$M \equal N$، آنگاه~$PM \equal PN$.
  \item اگر~$M \equal N$، آنگاه~$MQ \equal NQ$.
  \item اگر~$M \equal N$، آنگاه~$\lambd[x][M] \equal \lambd[x][N]$.
  \item $(\lambd[x][N])Q \equal \Subst{N}{Q}{x}$.
  \end{enumerate}
\end{defn}

سه قاعدهٔ نخست رابطه را به رابطه‌ای هم‌ارزی بدل می‌کنند؛ سه قاعدهٔ بعدی
آن را سازگار می‌کنند؛ و قاعدهٔ آخر تضمین می‌کند که رابطه، انقباض~$\beta$
را در بر می‌گیرد.

به‌طور غیرصوری، دو ترم هم‌ارز به~$\beta$ هستند اگر و تنها اگر بتوان یکی
را با صفر یا چند گام انقباض~$\beta$، یا «وارون» انقباض~$\beta$، به دیگری
تبدیل کرد. وارون انقباض~$\beta$ چنان تعریف می‌شود که~$M$ با وارون
انقباض~$\beta$ به~$N$ تبدیل شود اگر و تنها اگر~$N$ با انقباض~$\beta$
به~$M$ تبدیل شود.

علاوه بر قواعد بالا، رابطه را با قواعد بیشتری گسترش خواهیم داد و رابطهٔ
هم‌ارزی گسترش‌یافته را با~$\equal[X]$ نشان خواهیم داد، که در آن~$X$
قاعدهٔ گسترش‌دهنده است.

\end{document}
